\documentclass[11pt, a4paper]{article}

\usepackage[margin=2.4cm]{geometry}
\usepackage{amsmath, amssymb, amsthm}
\usepackage{mathtools}
\usepackage{bm}
\usepackage{enumitem}
\usepackage{booktabs}
\usepackage{hyperref}
\usepackage[numbers,sort&compress]{natbib}
\usepackage{microtype}
\usepackage{xcolor}
\usepackage{algorithm}
\usepackage{algorithmic}
\usepackage{appendix}

\definecolor{accent}{HTML}{2E5090}
\hypersetup{colorlinks=true, linkcolor=accent, citecolor=accent, urlcolor=accent}

\newcommand{\E}{\mathcal{E}}
\newcommand{\R}{\mathbb{R}}
\newcommand{\bx}{\bm{x}}
\newcommand{\by}{\bm{y}}
\newcommand{\bb}{\bm{b}}
\newcommand{\br}{\bm{r}}
\newcommand{\btheta}{\bm{\theta}}
\newcommand{\balpha}{\bm{\alpha}}
\newcommand{\bbeta}{\bm{\beta}}
\newcommand{\bc}{\bm{c}}
\newcommand{\bu}{\bm{u}}
\newcommand{\bv}{\bm{v}}
\newcommand{\bz}{\bm{z}}
\newcommand{\cG}{\mathcal{G}}
\newcommand{\cL}{\mathcal{L}}
\newcommand{\cS}{\mathcal{S}}
\newcommand{\cC}{\mathcal{C}}
\newcommand{\cB}{\mathcal{B}}
\newcommand{\cI}{\mathcal{I}}

\title{\bfseries Structure Learning as Precision Matching}
\author{}
\date{}

\begin{document}
\maketitle
\thispagestyle{empty}

%═══════════════════════════════════════════════════════════════
\begin{abstract}
%═══════════════════════════════════════════════════════════════
We introduce a method for selecting the graph topology of energy-based models at random initialisation, before any training.
The method scores each candidate graph by how well the \emph{directional structure} of its boundary precision matches the task: the Schur complement of the inference precision matrix, after marginalising internal degrees of freedom, is compared against the empirical task covariance via an A-optimal criterion.
Because each edge's Jacobian inherits the algebraic structure of its transform type---Fourier eigenvectors for convolutions, isotropic for dense layers, position-averaged for attention---the boundary Schur complement carries a directional signature that is determined by the transform family, not by random weight values.
The score is low when this directional signature aligns with the principal directions of the task covariance, providing a principled mechanism for selecting transform types: convolutions are preferred on data with structured frequency content, attention on data with positional structure.
Because the precision matrix decomposes additively over factor-graph edges, internal edges can only \emph{increase} the boundary Schur complement (in the PSD sense), making backward elimination monotonic and well-posed.
Dead-end edges contribute nothing to the Schur complement and are pruned for free; chain edges propagate precision from input to output; skip connections add direct boundary coupling.
The algorithm requires no training, no inference simulation, and no driving force---only the Jacobians at a single linearisation point.
\end{abstract}

%═══════════════════════════════════════════════════════════════
\section{Introduction}\label{sec:intro}
%═══════════════════════════════════════════════════════════════

Consider any model in which latent variables are adjusted to minimise a sum of local energies before parameters are updated: predictive coding networks~\citep{rao1999predictive,whittington2017approximation}, deep equilibrium models~\citep{bai2019deep}, Hopfield networks~\citep{ramsauer2021hopfield}, factor graphs in probabilistic graphical models~\citep{koller2009pgm}.
A factor graph $\cG$ connects variables $\bv_i \in \R^{d_i}$, partitioned into clamped data $\by$ and free latent states $\bx \in \R^D$, through parameterised transforms $f_1,\dots,f_M$ with parameters $\btheta$.
Local energies $\E_i$ penalise prediction discrepancies, giving total energy $E(\bx;\btheta,\cG,\by) = \sum_i \E_i$.
We address a question that precedes both inference and learning: \emph{which edges does this graph need?}

The energy decomposes over factors, so its Gauss--Newton Hessian decomposes additively over edges, giving an inference precision matrix $\bm{\Gamma}_\cG = \sum_e \bm{J}_e^\top \bm{\Lambda}_e \bm{J}_e$.
This is a classical object~\citep{lauritzen1996graphical,minka2001ep}.
What is new is extracting from it a score that measures \emph{how well the directional structure of the graph's boundary precision matches the task}---not how well inference performs with random weights, but whether the transform types chosen for each edge produce precision in directions that align with the data covariance.

The key object is the \textbf{boundary Schur complement}: the effective precision matrix on boundary-connected variables after marginalising internal degrees of freedom.
This is a property of the graph's Jacobian structure alone---no driving force, no inference simulation, no dependence on the random initial predictions.
Internal edges contribute to it only through the off-diagonal coupling they create between boundary-connected variables; dead-end edges contribute nothing.
Because the precision matrix is additive over edges and the Schur complement is monotone in the PSD order, removing an internal edge can only \emph{decrease} the boundary Schur complement, making backward elimination well-posed.

The directional structure of $\bm{\Gamma}_\cB^*$ is not random---it is inherited from the Jacobian algebra of the graph's transforms.
A convolutional edge imposes Fourier eigenvectors; a dense edge imposes isotropic random eigenvectors; an attention edge imposes position-averaged directions.
The score measures the alignment between these transform-determined directions and the principal variation directions of the task data.
This is the mechanism by which the method selects not just topology (which variables to connect) but transform type (how to connect them).

The boundary Schur complement connects two classical ideas.
In graphical models, the Schur complement of the precision matrix gives the marginal precision after variable elimination~\citep{lauritzen1996graphical}.
In Bayesian model reduction~\citep{friston2018bmr}, a full model is inverted once and reduced models are scored analytically.
Our method applies the same \emph{invert-once-reduce-many} logic, but scores by boundary precision rather than variational free energy, and operates on graph edges rather than parameter priors.

%═══════════════════════════════════════════════════════════════
\section{The Boundary Schur Complement}\label{sec:score}
%═══════════════════════════════════════════════════════════════

\subsection{The inference precision matrix}\label{sec:construction}

Linearising each transform $f_e$ around the current state $\bx_0$, the Gauss--Newton approximation to the energy Hessian is (Appendix~\ref{app:hessian}):
\begin{equation}\label{eq:hessian}
    \nabla^2_{\bx} E \;\approx\; \bm{\Gamma}_\cG \;=\; \sum_e \bm{J}_e^\top \bm{\Lambda}_e\,\bm{J}_e,
\end{equation}
where $\bm{J}_e \in \R^{d_{\mathrm{tgt}(e)}\times D}$ is the Jacobian of the prediction residual with respect to the free states, and $\bm{\Lambda}_e \succeq 0$ is the curvature of the local energy at edge $e$.
For linear transforms this is exact; for nonlinear transforms it is accurate when prediction errors are small or transforms are near-linear (Appendix~\ref{app:hessian}).
Each edge contributes a PSD term $\bm{J}_e^\top\bm{\Lambda}_e\bm{J}_e$ to the sum.

\subsection{Partitioning the state space}\label{sec:partition}

Partition the $D$ free state dimensions into two groups:
\begin{itemize}[nosep,leftmargin=2em]
    \item $B$-dimensions: those belonging to variables touched by at least one boundary edge (an edge with a clamped endpoint).
    These are the variables that directly interface with observed data.
    \item $I$-dimensions: those belonging to variables touched only by internal edges.
    These are purely latent variables with no direct data connection.
\end{itemize}
Under this partition, the precision matrix has block structure:
\begin{equation}\label{eq:block}
    \bm{\Gamma}_S = \begin{bmatrix} \bm{\Gamma}_{BB} & \bm{\Gamma}_{BI} \\ \bm{\Gamma}_{IB} & \bm{\Gamma}_{II} \end{bmatrix} + \varepsilon\bm{I},
\end{equation}
where $S$ is the active edge set and $\varepsilon > 0$ provides regularisation.

\subsection{The Schur complement as boundary precision}\label{sec:schur}

The \textbf{boundary Schur complement} is the effective precision on $B$-dimensions after marginalising $I$-dimensions:
\begin{equation}\label{eq:schur}
    \bm{\Gamma}_\cB^*(S) \;=\; \bm{\Gamma}_{BB} - \bm{\Gamma}_{BI}\,\bm{\Gamma}_{II}^{-1}\,\bm{\Gamma}_{IB}.
\end{equation}
In the Gaussian model $p(\bx) \propto \exp(-\frac{1}{2}\bx^\top\bm{\Gamma}_S\bx)$, this is exactly the marginal precision of $\bx_B$~\citep{lauritzen1996graphical}.
It measures how much information flows from input boundary to output boundary through the graph's internal structure.

The score for edge set $S$ is:
\begin{equation}\label{eq:score_def}
    J(S) = \mathrm{tr}\!\Big((\bm{\Gamma}_\cB^*(S))^{-1}\,\bm{\Sigma}_{\mathrm{task}}\Big),
\end{equation}
where $\bm{\Sigma}_{\mathrm{task}} = \frac{1}{B}\sum_n \bb_{\cB,n}\bb_{\cB,n}^\top$ is the empirical covariance of boundary residuals over the probe batch, and $\bb_{\cB,n}$ stacks the curvature-weighted prediction errors at boundary edges for sample $n$.
This is an A-optimal design criterion~\citep{pukelsheim2006optimal}: the graph is scored by expected residual variance in the boundary subspace, weighted by the task distribution.
Lower is better---more boundary precision means less expected residual.

The trace has a directional interpretation.
The eigenvectors of $\bm{\Gamma}_\cB^*$ are inherited from the Jacobian algebra of the graph's transforms: Fourier basis for convolutions, isotropic random for dense layers, position-averaged for attention (Section~\ref{sec:structure}).
The eigenvectors of $\bm{\Sigma}_{\mathrm{task}}$ reflect the principal variation directions of the data at the boundary.
The score $J(S)$ is small when these two sets of eigenvectors are well-aligned---when the graph's transforms produce precision in exactly the directions where the task has the most variance.
This is the mechanism by which the score selects transform types: convolutions are preferred on data with structured frequency content (natural images) because their Fourier precision aligns with the data's frequency covariance; attention is preferred on data with positional structure because its position-averaged precision aligns with positional variation.

Three properties make this score suitable for structure learning.

\paragraph{No driving force.}
The score depends only on the precision matrix $\bm{\Gamma}_S$, not on any inference simulation or initial prediction quality.
At random initialisation, individual edge predictions are random noise, but the \emph{structural} coupling they induce in $\bm{\Gamma}_\cB^*$ is deterministic given the Jacobians.
The score measures structural capacity, not random-init performance.

\paragraph{Monotonicity.}
Adding a PSD term to $\bm{\Gamma}_S$ can only increase $\bm{\Gamma}_\cB^*$ in the PSD order (Appendix~\ref{app:monotone}).
Therefore removing an internal edge can only decrease $\bm{\Gamma}_\cB^*$, which can only increase $J(S)$.
Backward elimination is well-posed: the score is monotonically non-decreasing as internal edges are removed.

\paragraph{Dead-end insensitivity.}
An edge $h_i \to d_j$ where $d_j$ is touched only by internal edges contributes PSD mass only in the $I$-block of $\bm{\Gamma}_S$.
It adds to $\bm{\Gamma}_{II}$ but not to $\bm{\Gamma}_{BI}$ or $\bm{\Gamma}_{BB}$.
Its contribution to the Schur complement is exactly zero---it cannot be ``seen'' from the boundary.
Such edges have zero leverage and are pruned for free.

\subsection{Coverage and conditioning}\label{sec:coverage}

The boundary Schur complement $\bm{\Gamma}_\cB^*$ has eigendecomposition $\bm{\Gamma}_\cB^* = \bm{V}\bm{\Omega}\bm{V}^\top$ with eigenvalues $\omega_1 \geq \cdots \geq \omega_{D_B} \geq 0$.
The eigenvectors $\bm{V}$ are inherited from the Jacobian structure of the graph's transforms, propagated through the Schur complement; the eigenvalues $\omega_i$ reflect how much precision each direction carries.
The score decomposes as:
\begin{equation}\label{eq:decomp}
    J(S) = \sum_{i=1}^{D_B} \frac{\sigma_i}{\omega_i + \varepsilon},
\end{equation}
where $\sigma_i = \bv_i^\top\bm{\Sigma}_{\mathrm{task}}\bv_i$ is the task variance along eigendirection $\bv_i$.
This separates into two failure modes:
\begin{itemize}[nosep,leftmargin=2em]
    \item \textbf{Coverage gap}: directions with $\omega_i \approx 0$ contribute $\sigma_i/\varepsilon$---large, regularisation-dominated terms.
    These are boundary directions that no internal path connects, structurally unreachable regardless of inference budget.
    \item \textbf{Conditioning penalty}: directions with $0 < \omega_i \ll 1/\eta T$ contribute terms that are resolved too slowly.
    These are reachable but poorly conditioned---the inference-side analogue of vanishing gradients through deep chains.
\end{itemize}
A good architecture minimises both: full coverage (every task-relevant boundary direction is reachable through internal structure) and good conditioning (reachable directions have large eigenvalues in $\bm{\Gamma}_\cB^*$).

\subsection{Why finite $T$ matters}\label{sec:whyT}

At convergence ($T\to\infty$, equivalently $\varepsilon\to 0$), the score depends only on whether $\omega_i > 0$---the rank of $\bm{\Gamma}_\cB^*$.
A 100-layer chain and a single skip connection score identically if both provide full-rank boundary coupling.
There is no basis for preferring one over the other.

The conditioning penalty restores discrimination.
Setting $\varepsilon = 1/\eta T$ (matching the spectral gate of $T$-step gradient descent), the score penalises boundary directions whose precision $\omega_i$ falls below $1/\eta T$.
A deep chain has exponentially spread eigenvalues in $\bm{\Gamma}_\cB^*$---most below the gate.
A skip connection has $O(1)$ eigenvalues---above the gate.
The inference budget $T$ is the mechanism that makes structure learning non-trivial.

\subsection{Relation to model evidence}\label{sec:evidence}

The Laplace-approximated log marginal likelihood includes an Occam factor $\frac{1}{2}\log\det(\bm{\Gamma}_\cG)$ penalising \emph{large} eigenvalues (over-committed capacity).
Our score penalises \emph{small} eigenvalues of the boundary Schur complement (computationally inaccessible capacity).
These are complementary: model evidence asks ``does this graph overfit?''; our score asks ``can this graph resolve the boundary task?''
Bayesian model reduction~\citep{friston2018bmr,friston2016group} scores reduced models by the change in variational free energy; we score by the change in boundary precision.
Both share the \emph{invert-once-reduce-many} structure (Section~\ref{sec:woodbury}).

%═══════════════════════════════════════════════════════════════
\section{Backward Elimination}\label{sec:search}
%═══════════════════════════════════════════════════════════════

\subsection{Algorithm}\label{sec:algorithm}

Given per-edge Jacobians $\bm{J}_e$ (computed from a single forward pass on the supergraph) and boundary residuals $\bb_{\cB,n}$ (from the probe batch), the score of any edge set $S$ is computable in $O(D^3)$ via the Schur complement~\eqref{eq:schur} and the trace~\eqref{eq:score_def}.

\begin{algorithm}[t]
\caption{Backward Elimination via Boundary Schur Complement}
\label{alg:exact}
\begin{algorithmic}[1]
\REQUIRE Supergraph with $M$ edges, probe batch $\{(\bx_n,\by_n)\}_{n=1}^N$, gate $\varepsilon = 1/\eta T$
\STATE One forward pass: compute $\bm{J}_e$, $\bm{\Lambda}_e$ for all edges; $\bb_{\cB,n}$ for all samples
\STATE Classify edges: $\cB$ (boundary, fixed), $\cI$ (internal, prunable)
\STATE Compute $\bm{\Sigma}_{\mathrm{task}} = \frac{1}{N}\sum_n\bb_{\cB,n}\bb_{\cB,n}^\top$
\STATE $S \leftarrow \{1,\dots,M\}$; \; $J_{\mathrm{curr}} \leftarrow J(S)$
\REPEAT
    \FOR{each $e \in S \cap \cI$}
        \STATE $J_e \leftarrow J(S\setminus\{e\})$ \hfill\COMMENT{$O(D^3)$ via Schur complement}
    \ENDFOR
    \STATE $e^* \leftarrow \arg\min_{e \in S \cap \cI} J_e$
    \IF{$J_{e^*} > J_{\mathrm{curr}} + \delta$}
        \STATE \textbf{break} \hfill\COMMENT{every removal exceeds tolerance}
    \ENDIF
    \STATE $S \leftarrow S \setminus \{e^*\}$; \; $J_{\mathrm{curr}} \leftarrow J_{e^*}$
\UNTIL{$S \cap \cI = \emptyset$}
\RETURN Edge set $S$, Schur complement $\bm{\Gamma}_\cB^*(S)$, score decomposition
\end{algorithmic}
\end{algorithm}

\paragraph{Boundary edges are never pruned.}
They define the task interface---removing one changes which data the graph can see.

\paragraph{Stopping criterion.}
By monotonicity, $J(S\setminus\{e\}) \geq J(S)$ for all internal $e$.
Dead-end edges give $J(S\setminus\{e\}) = J(S)$ exactly (zero contribution to Schur complement) and are removed for free.
The algorithm stops when the cheapest removal exceeds tolerance $\delta$---every remaining edge contributes meaningful boundary precision.

\paragraph{Cost.}
Each step evaluates $|S\cap\cI|$ candidates at $O(D^3)$ each (or $O(D_B^3)$ if only the Schur complement is recomputed).
Total: $O(M^2 D_B^3)$.
For $D_B \lesssim 200$ and $M \lesssim 100$, this completes in seconds.

\subsection{Scalable approximation via Woodbury}\label{sec:woodbury}

For large $D$, maintain $\bm{\Gamma}_S^{-1}$ and update incrementally.
The Schur complement can be read from $\bm{\Gamma}_S^{-1}$: $(\bm{\Gamma}_\cB^*)^{-1} = [\bm{\Gamma}_S^{-1}]_{BB}$, the $B$-block of the full inverse~\citep{lauritzen1996graphical}.
Removing edge $e$ updates $\bm{\Gamma}_S^{-1}$ via:
\begin{equation}\label{eq:woodbury}
    (\bm{\Gamma}_S - \bm{J}_e^\top\bm{\Lambda}_e\bm{J}_e)^{-1} = \bm{\Gamma}_S^{-1} + \bm{\Gamma}_S^{-1}\bm{J}_e^\top(\bm{\Lambda}_e^{-1} - \bm{J}_e\bm{\Gamma}_S^{-1}\bm{J}_e^\top)^{-1}\bm{J}_e\bm{\Gamma}_S^{-1},
\end{equation}
reducing per-candidate cost from $O(D^3)$ to $O(D^2 d_{\mathrm{tgt}})$.
The updated score is read from the $B$-block of the downdated inverse.
This is the structural analogue of BMR: invert once, score every reduction analytically.

%═══════════════════════════════════════════════════════════════
\section{How Structure Enters the Schur Complement}\label{sec:structure}
%═══════════════════════════════════════════════════════════════

Each edge contributes a PSD term to $\bm{\Gamma}_S$, which enters the Schur complement $\bm{\Gamma}_\cB^*$ through the block structure~\eqref{eq:block}.
We trace how topology and transform families affect boundary precision at random initialisation.

\paragraph{Chain edges.}
A chain $\bv_0 \to h_1 \to \cdots \to h_L \to \bv_L$ where $\bv_0$ and $\bv_L$ are boundary-connected creates off-diagonal coupling $\bm{\Gamma}_{BI}$ between boundary and internal dimensions.
The Schur complement $\bm{\Gamma}_{BB} - \bm{\Gamma}_{BI}\bm{\Gamma}_{II}^{-1}\bm{\Gamma}_{IB}$ propagates precision from input to output through the chain.
Under random initialisation, the effective coupling decays exponentially with depth (singular values of the composed Jacobian spread as $e^{cL}$), so the boundary precision from a deep chain has exponentially spread eigenvalues---most below the spectral gate.

\paragraph{Skip connections.}
A skip edge $\bv_i \to \bv_j$ between two boundary-connected variables adds directly to $\bm{\Gamma}_{BB}$---no marginalisation needed.
Its contribution has $O(1)$ eigenvalues at random init.
In deep graphs, skip connections dominate the Schur complement because the chain's contribution is exponentially attenuated.

\paragraph{Dead-end edges.}
An edge $h_i \to d_j$ where $d_j$ has no boundary connection contributes only to $\bm{\Gamma}_{II}$.
It does not appear in $\bm{\Gamma}_{BI}$ or $\bm{\Gamma}_{BB}$.
Its contribution to the Schur complement is exactly zero.

\paragraph{Parallel edges.}
$K$ parallel edges contribute $K$ independent PSD terms.
The first provides new rank in $\bm{\Gamma}_\cB^*$ (large leverage); additional ones add eigenvalue mass on existing directions (decreasing marginal leverage).
The algorithm naturally thins parallel edges to the minimum set.

\paragraph{Dense transforms.}
$f(\bv) = \bm{W}\bv$ with i.i.d.\ entries gives a Wishart contribution with isotropic eigenvectors (Haar-distributed singular vectors).
No direction is preferred; the boundary precision aligns with the task no better than chance.
Dense transforms provide capacity without directional bias---a neutral baseline against which structured transforms are measured.

\paragraph{Convolutional transforms.}
$f(\bv) = \bm{W}*\bv$ gives a block-circulant Jacobian diagonalised by the DFT, regardless of filter values.
The Schur complement inherits Fourier eigenvectors: boundary precision decomposes by spatial frequency.
On natural images, whose task covariance $\bm{\Sigma}_{\mathrm{task}}$ concentrates variance at low spatial frequencies, this alignment produces a lower score than isotropic alternatives---the mechanism by which the algorithm prefers convolutions on images.

\paragraph{Attention.}
At random init, the Jacobian approximates mean-pooling with a random value projection, contributing position-averaged precision to $\bm{\Gamma}_\cB^*$.
On sequential data where $\bm{\Sigma}_{\mathrm{task}}$ has positional structure, this alignment is favourable.

\paragraph{Recurrence.}
Jacobians involve matrix products with singular values decaying with lag.
The Schur complement from a recurrent path attenuates long-range boundary coupling.
Gated variants preserve singular values, maintaining boundary precision over longer horizons.

%═══════════════════════════════════════════════════════════════
\section{Related Work}\label{sec:related}
%═══════════════════════════════════════════════════════════════

\paragraph{Bayesian model reduction.}
BMR~\citep{friston2018bmr,friston2016group} fits a full model, then analytically scores reduced models differing only in priors.
Our ridge approximation (Section~\ref{sec:woodbury}) inherits this structure but operates on graph edges and scores by boundary precision rather than variational free energy.
BMR penalises large eigenvalues (statistical over-commitment); we penalise small eigenvalues of the Schur complement (computational inaccessibility).
BMR has been applied to neural network pruning~\citep{beckers2024bmrs}; we prune edges---the natural structural unit for locally factored energies.

\paragraph{Structure learning in graphical models.}
The Schur complement of the precision matrix as marginal precision is foundational to Gaussian graphical models~\citep{lauritzen1996graphical,koller2009pgm}.
Score-based structure learning selects graphs by penalised likelihood or independence tests.
Our approach uses the same object at random initialisation with heterogeneous neural transforms and a finite-budget spectral gate.

\paragraph{Training-free NAS.}
TE-NAS~\citep{chen2021tenas} and NASI~\citep{shu2022nasi} use NTK-derived quantities~\citep{jacot2018neural} to rank architectures, but NTK metrics conflate component spectra through end-to-end composition~\citep{mok2022demystifying,zhou2024vintk}.
The Schur complement decomposes over local factors and marginalises internal structure, avoiding this conflation.

\paragraph{Optimal experimental design.}
The score~\eqref{eq:score_def} is A-optimal: it minimises expected posterior variance at the boundary under the task distribution~\citep{pukelsheim2006optimal}.
Greedy backward elimination for A-optimal design has approximate submodularity guarantees~\citep{krause2008near}, which may provide non-trivial bounds on the quality of the reduced graph.

\paragraph{Sparse approximation.}
Forward structure search---adding edges greedily to increase boundary precision---is matching pursuit~\citep{mallat1993matching} in Jacobian space.
Compressed sensing theory~\citep{tropp2004greed} provides guarantees under dictionary coherence conditions.

%═══════════════════════════════════════════════════════════════
\section{Discussion}\label{sec:discussion}
%═══════════════════════════════════════════════════════════════

We have presented a structure-learning method for energy-based models that scores graph topology by how well the directional structure of its boundary precision matches the task.

The method rests on three observations.
First, each transform type imposes a specific eigenvector basis on its contribution to the precision matrix---Fourier for convolutions, isotropic for dense, position-averaged for attention---and this basis is an algebraic property of the transform class, not an artifact of random weights.
The Schur complement propagates this directional structure to the boundary, where the A-optimal score measures its alignment with the task covariance.
This is the mechanism that drives architecture selection: convolutions are preferred on images because their Fourier precision matches the frequency structure of natural image data.
Second, the energy's local factorisation makes the precision matrix additive over edges, and the Schur complement inherits monotonicity: adding an internal edge can only increase boundary precision.
Third, the spectral gate $\varepsilon = 1/\eta T$ connects the inference budget to the score: at convergence all connected graphs look alike, but at finite $T$ the conditioning of the Schur complement discriminates between them.

The algorithm is simple: backward elimination, removing the internal edge whose removal least increases the boundary score.
Dead-end edges have zero contribution and are removed for free.
Chain edges propagate precision through internal variables and survive until their exponentially attenuated coupling falls below the gate.
Skip connections add direct boundary coupling and survive longest.

The framework makes concrete, mechanistically grounded predictions.
Dead-end edges should be pruned first (zero Schur complement contribution).
Chain edges should be pruned before skip connections in deep graphs (exponentially attenuated boundary coupling vs.\ $O(1)$ direct coupling).
Convolutional edges should be preferred over dense edges on natural images (Fourier precision aligns with frequency-structured task covariance).
Attention edges should be preferred on sequential data (position-averaged precision aligns with positional task covariance).
The pruning order should be $T$-sensitive: short budgets favour direct paths, long budgets tolerate depth (the spectral gate determines which conditioning levels are acceptable).
Each prediction follows from a specific algebraic property of the transforms, not from a fitted correlation.

%═══════════════════════════════════════════════════════════════

\bibliographystyle{plainnat}
\begin{thebibliography}{17}

\bibitem[Bai et~al.(2019)]{bai2019deep}
S.~Bai, J.~Z.~Kolter, and V.~Koltun.
\newblock Deep equilibrium models.
\newblock In \emph{NeurIPS}, 2019.

\bibitem[Beckers et~al.(2024)]{beckers2024bmrs}
S.~Beckers, J.~Jantre, and others.
\newblock {BMRS}: {B}ayesian model reduction for structured pruning.
\newblock \emph{arXiv:2406.01345}, 2024.

\bibitem[Chen et~al.(2021)]{chen2021tenas}
W.~Chen, X.~Gong, and Z.~Wang.
\newblock Neural architecture search on {ImageNet} in four {GPU} hours.
\newblock In \emph{ICLR}, 2021.

\bibitem[Friston et~al.(2016)]{friston2016group}
K.~J.~Friston et~al.
\newblock {B}ayesian model reduction and empirical {B}ayes for group ({DCM}) studies.
\newblock \emph{NeuroImage}, 128:413--431, 2016.

\bibitem[Friston et~al.(2018)]{friston2018bmr}
K.~Friston, T.~Parr, and P.~Zeidman.
\newblock {B}ayesian model reduction.
\newblock \emph{arXiv:1805.07092}, 2018.

\bibitem[Gray(2006)]{gray2006toeplitz}
R.~M.~Gray.
\newblock \emph{Toeplitz and Circulant Matrices: A Review}.
\newblock Now Publishers, 2006.

\bibitem[Jacot et~al.(2018)]{jacot2018neural}
A.~Jacot, F.~Gabriel, and C.~Hongler.
\newblock Neural tangent kernel: Convergence and generalization in neural networks.
\newblock In \emph{NeurIPS}, 2018.

\bibitem[Koller and Friedman(2009)]{koller2009pgm}
D.~Koller and N.~Friedman.
\newblock \emph{Probabilistic Graphical Models: Principles and Techniques}.
\newblock MIT Press, 2009.

\bibitem[Krause et~al.(2008)]{krause2008near}
A.~Krause, A.~Singh, and C.~Guestrin.
\newblock Near-optimal sensor placements in {Gaussian} processes: Theory, efficient algorithms and empirical studies.
\newblock \emph{JMLR}, 9:235--284, 2008.

\bibitem[Lauritzen(1996)]{lauritzen1996graphical}
S.~L.~Lauritzen.
\newblock \emph{Graphical Models}.
\newblock Oxford University Press, 1996.

\bibitem[Mallat and Zhang(1993)]{mallat1993matching}
S.~Mallat and Z.~Zhang.
\newblock Matching pursuits with time-frequency dictionaries.
\newblock \emph{IEEE TSP}, 41(12):3397--3415, 1993.

\bibitem[Marchenko and Pastur(1967)]{marchenko1967distribution}
V.~A.~Marchenko and L.~A.~Pastur.
\newblock Distribution of eigenvalues for some sets of random matrices.
\newblock \emph{Matematicheskii Sbornik}, 114(4):507--536, 1967.

\bibitem[Minka(2001)]{minka2001ep}
T.~P.~Minka.
\newblock Expectation propagation for approximate {Bayesian} inference.
\newblock In \emph{UAI}, 2001.

\bibitem[Mok et~al.(2022)]{mok2022demystifying}
J.~Mok, B.~Na, H.~Choe, and S.~Yoon.
\newblock Demystifying the neural tangent kernel from a practical perspective.
\newblock In \emph{CVPR}, 2022.

\bibitem[Pukelsheim(2006)]{pukelsheim2006optimal}
F.~Pukelsheim.
\newblock \emph{Optimal Design of Experiments}.
\newblock SIAM, 2006.

\bibitem[Ramsauer et~al.(2021)]{ramsauer2021hopfield}
H.~Ramsauer et~al.
\newblock Hopfield networks is all you need.
\newblock In \emph{ICLR}, 2021.

\bibitem[Rao and Ballard(1999)]{rao1999predictive}
R.~P.~N.~Rao and D.~H.~Ballard.
\newblock Predictive coding in the visual cortex.
\newblock \emph{Nature Neuroscience}, 2(1):79--87, 1999.

\bibitem[Shu et~al.(2022)]{shu2022nasi}
Y.~Shu, Z.~Cai, Z.~Dai, T.~Ooi, and B.~K.~H.~Low.
\newblock {NASI}: Label- and data-agnostic neural architecture search at initialization.
\newblock In \emph{ICLR}, 2022.

\bibitem[Tropp(2004)]{tropp2004greed}
J.~A.~Tropp.
\newblock Greed is good: Algorithmic results for sparse approximation.
\newblock \emph{IEEE Trans.\ Inform.\ Theory}, 50(10):2231--2242, 2004.

\bibitem[Whittington and Bogacz(2017)]{whittington2017approximation}
J.~C.~R.~Whittington and R.~Bogacz.
\newblock An approximation of the error backpropagation algorithm in a predictive coding network with local {Hebbian} synaptic plasticity.
\newblock \emph{Neural Computation}, 29(5):1229--1262, 2017.

\bibitem[Winn and Bishop(2005)]{winn2005variational}
J.~Winn and C.~M.~Bishop.
\newblock Variational message passing.
\newblock \emph{JMLR}, 6:661--694, 2005.

\bibitem[Zhou et~al.(2024)]{zhou2024vintk}
W.~Zhou, J.~Huang, and others.
\newblock When training-free {NAS} meets vision transformers: A neural tangent kernel perspective.
\newblock In \emph{DAC}, 2024.

\end{thebibliography}

%═══════════════════════════════════════════════════════════════
\newpage
\begin{appendices}

\section{Gauss--Newton Decomposition}\label{app:hessian}

For a single energy term $\E_i(\bv_i,\hat{\bv}_i)$ with prediction $\hat{\bv}_i = f_e(\bv_{\mathrm{src}(e)})$, the Hessian with respect to free variables $\bz$ is
\begin{equation}
    \nabla^2_{\bz}\E_i = \bm{J}_e^{(\bz)\top}\bm{\Lambda}_e\bm{J}_e^{(\bz)} + \sum_j(\nabla_{\hat{v}_j}\E_i)\nabla^2_{\bz}[f_e]_j.
\end{equation}
The second term is zero for linear transforms and negligible when prediction errors are small.
Summing over edges gives~\eqref{eq:hessian}.

\section{Monotonicity of the Schur Complement}\label{app:monotone}

Let $\bm{\Gamma}$ have the block structure~\eqref{eq:block} and let $\bm{\Delta} \succeq 0$ be a PSD matrix supported on the $I$-block (i.e., $\bm{\Delta}_{BB} = 0$, $\bm{\Delta}_{BI} = 0$).
Then:
\begin{align}
    \bm{\Gamma}_\cB^*(\bm{\Gamma}) &= \bm{\Gamma}_{BB} - \bm{\Gamma}_{BI}\bm{\Gamma}_{II}^{-1}\bm{\Gamma}_{IB}, \\
    \bm{\Gamma}_\cB^*(\bm{\Gamma}+\bm{\Delta}) &= \bm{\Gamma}_{BB} - \bm{\Gamma}_{BI}(\bm{\Gamma}_{II}+\bm{\Delta}_{II})^{-1}\bm{\Gamma}_{IB}.
\end{align}
Since $\bm{\Delta}_{II} \succeq 0$, we have $(\bm{\Gamma}_{II}+\bm{\Delta}_{II})^{-1} \preceq \bm{\Gamma}_{II}^{-1}$, so $\bm{\Gamma}_{BI}(\bm{\Gamma}_{II}+\bm{\Delta}_{II})^{-1}\bm{\Gamma}_{IB} \preceq \bm{\Gamma}_{BI}\bm{\Gamma}_{II}^{-1}\bm{\Gamma}_{IB}$, and therefore $\bm{\Gamma}_\cB^*(\bm{\Gamma}+\bm{\Delta}) \succeq \bm{\Gamma}_\cB^*(\bm{\Gamma})$.

For a general internal edge whose PSD contribution spans both $B$ and $I$ blocks (because its Jacobian touches a boundary-connected variable), the result extends: the full Schur complement $\bm{\Gamma}_{BB} + \bm{\Delta}_{BB} - (\bm{\Gamma}_{BI}+\bm{\Delta}_{BI})(\bm{\Gamma}_{II}+\bm{\Delta}_{II})^{-1}(\bm{\Gamma}_{IB}+\bm{\Delta}_{IB})$ is at least $\bm{\Gamma}_\cB^*(\bm{\Gamma})$, by the generalised Schur complement monotonicity theorem for PSD updates to the full matrix~\citep{lauritzen1996graphical}.

\end{appendices}

\end{document}